\section{Erd\H{o}s Problem \#301}

\subsection*{1) ROUND-2 OBJECTIVE}

I pursue path (A): prove a genuinely stronger upper bound for $f(N)$.

Round~1 reconstructed the known bound
\[
f(N)\le \left(\frac{25}{28}+o(1)\right)N
\]
from the disjoint blocks $\{2a,3a,4a,6a,12a\}$. The most promising Round-2 direction is to enlarge the local forbidden block while keeping the blocks pairwise disjoint. This is exactly what the new argument does.

\subsection*{2) Round-1 FOUNDATION USED}

I use the following Round-1 facts.

\begin{itemize}
\item The definition of $f(N)$ as the largest size of a set $A\subseteq\{1,\ldots,N\}$ containing no distinct solution to
\[
\frac1a=\frac1{b_1}+\cdots+\frac1{b_k}.
\]
\item The trivial lower bound $f(N)\ge \lfloor N/2\rfloor$, from $A$ equal to the integers in $(N/2,N]$.
\item The Round-1 block argument giving
\[
f(N)\le \left(\frac{25}{28}+o(1)\right)N.
\]
\item The exact values of $f(N)$ for $N\le 15$.
\end{itemize}

I do not re-prove the Round-1 $25/28$ bound; I only use it as the baseline to improve upon.

\subsection*{3) NEW INSIGHT / TOOL (ROUND-2)}

The new block is
\[
T_a:=\{2a,3a,4a,6a,8a,12a,24a\}\cap[1,N].
\]

I choose $a$ from the set
\[
\mathcal{I}:=\{16^b 9^c d:\ b,c\ge 0,\ (d,6)=1\}.
\]
Equivalently,
\[
v_2(a)\equiv 0\ (\mathrm{mod}\ 4),
\qquad
v_3(a)\equiv 0\ (\mathrm{mod}\ 2).
\]

The residue pairs modulo $(4,2)$ of the multipliers
\[
2,3,4,6,8,12,24
\]
are all distinct, so these larger blocks are still pairwise disjoint. The local hypergraph on this block is stronger than in Round~1 and forces more omissions.

\subsection*{4) ATTACK PLAN (ROUND-2)}

The exact Round-1 gap was quantitative: the local block $\{2a,3a,4a,6a,12a\}$ forced only a density loss of $3/28$.

To improve this I will:
\begin{itemize}
\item prove pairwise disjointness of the larger blocks $T_a$ for $a\in\mathcal I$;
\item prove local omission bounds:
\[
\begin{array}{ll}
a\le N/24 & \Longrightarrow \mbox{at least $3$ omissions from }T_a,\\
N/24<a\le N/12 & \Longrightarrow \mbox{at least $2$ omissions from }T_a,\\
N/12<a\le N/6 & \Longrightarrow \mbox{at least $1$ omission from }T_a;
\end{array}
\]
\item count how many such $a$ exist.
\end{itemize}

This directly overcomes the Round-1 obstacle by extracting more omissions from each disjoint block.

\subsection*{5) WORK (ROUND-2)}

\medskip
\noindent\textbf{Step 1: the new block and its identities.}

Write
\[
U_a:=\{2a,3a,4a,6a,8a,12a\}.
\]
Inside $T_a$ we have the exact identities
\[
\frac1{2a}=\frac1{3a}+\frac1{6a},
\qquad
\frac1{3a}=\frac1{4a}+\frac1{12a},
\qquad
\frac1{4a}=\frac1{6a}+\frac1{12a},
\]
and also
\[
\frac1{6a}=\frac1{8a}+\frac1{24a},
\qquad
\frac1{8a}=\frac1{12a}+\frac1{24a},
\qquad
\frac1{2a}=\frac1{3a}+\frac1{8a}+\frac1{24a}.
\]

\medskip
\noindent\textbf{Step 2: disjointness.}

\medskip
\noindent\textbf{Lemma 1.}
As $a$ ranges over $\mathcal I$, the sets $T_a$ are pairwise disjoint.

\medskip
\noindent\textbf{Proof.}
For $a\in\mathcal I$ we have
\[
v_2(a)\equiv 0\ (\mathrm{mod}\ 4),\qquad v_3(a)\equiv 0\ (\mathrm{mod}\ 2).
\]
If $x\in T_a$, then $x=ca$ for some
\[
c\in\{2,3,4,6,8,12,24\}.
\]
Therefore
\[
\bigl(v_2(x)\bmod 4,\ v_3(x)\bmod 2\bigr)
=
\bigl(v_2(c)\bmod 4,\ v_3(c)\bmod 2\bigr).
\]
The seven multipliers give the seven distinct residue pairs
\[
(1,0),\ (0,1),\ (2,0),\ (1,1),\ (3,0),\ (2,1),\ (3,1).
\]
Hence the residue pair determines the multiplier $c$, and then $a=x/c$ is uniquely determined. So no integer can belong to two different blocks $T_a$.

\medskip
\noindent\textbf{Step 3: counting the admissible $a$.}

Let
\[
I(X):=\#\{a\le X:\ a\in\mathcal I\}.
\]

\medskip
\noindent\textbf{Lemma 2.}
As $X\to\infty$,
\[
I(X)=\frac25 X+o(X).
\]

\medskip
\noindent\textbf{Proof.}
Write every $a\in\mathcal I$ uniquely as $a=16^b9^cd$ with $(d,6)=1$. If
\[
C(Y):=\#\{d\le Y:\ (d,6)=1\},
\]
then
\[
C(Y)=\frac13 Y+O(1).
\]
Therefore
\[
I(X)=\sum_{b,c\ge 0} C\!\left(\frac{X}{16^b9^c}\right)
=\sum_{b,c\ge 0}\left(\frac13\cdot \frac{X}{16^b9^c}+O(1)\right).
\]
Only $O((\log X)^2)$ pairs $(b,c)$ satisfy $16^b9^c\le X$, so the total error is $o(X)$. The main term is
\[
\frac{X}{3}\left(\sum_{b\ge 0}16^{-b}\right)\left(\sum_{c\ge 0}9^{-c}\right)
=
\frac{X}{3}\cdot \frac{16}{15}\cdot \frac98
=
\frac25 X.
\]
Hence
\[
I(X)=\frac25 X+o(X).
\]

\medskip
\noindent\textbf{Step 4: local omission lemmas.}

\medskip
\noindent\textbf{Lemma 3.}
If $a\le N/12$, then any admissible set $A\subseteq \{1,\ldots,N\}$ omits at least $2$ elements of $U_a$.

\medskip
\noindent\textbf{Proof.}
Assume for contradiction that $A$ omits at most one element of $U_a$. Then exactly one of the following happens.

\begin{itemize}
\item If the missing element is $2a$, $6a$, or $8a$, then $3a,4a,12a\in A$, and
\[
\frac1{3a}=\frac1{4a}+\frac1{12a}
\]
is a forbidden relation.
\item If the missing element is $3a$, then $4a,6a,12a\in A$, and
\[
\frac1{4a}=\frac1{6a}+\frac1{12a}
\]
is forbidden.
\item If the missing element is $4a$ or $12a$, then $2a,3a,6a\in A$, and
\[
\frac1{2a}=\frac1{3a}+\frac1{6a}
\]
is forbidden.
\end{itemize}

All cases contradict admissibility, so at least two elements of $U_a$ are omitted.

\medskip
\noindent\textbf{Lemma 4.}
If $a\le N/24$, then any admissible set $A\subseteq \{1,\ldots,N\}$ omits at least $3$ elements of $T_a$.

\medskip
\noindent\textbf{Proof.}
Suppose instead that at most two elements of $T_a$ are omitted.

If $24a\notin A$, then among the six elements of $U_a$ at most one is missing, contradicting Lemma~3.

So $24a\in A$. Then exactly two elements of $U_a$ are missing.

Now consider whether $8a$ lies in $A$.

\begin{itemize}
\item If $8a\notin A$, then four of the five elements
\[
2a,3a,4a,6a,12a
\]
belong to $A$. Any such $4$-subset contains one of the forbidden triples
\[
\{2a,3a,6a\},\qquad \{3a,4a,12a\},\qquad \{4a,6a,12a\},
\]
giving a contradiction.
\item If $8a\in A$, then to avoid
\[
\frac1{6a}=\frac1{8a}+\frac1{24a}
\]
we must have $6a\notin A$, and to avoid
\[
\frac1{8a}=\frac1{12a}+\frac1{24a}
\]
we must also have $12a\notin A$.
Since only two elements of $U_a$ are missing, this forces
\[
2a,3a,4a,8a,24a\in A.
\]
But then
\[
\frac1{2a}=\frac1{3a}+\frac1{8a}+\frac1{24a}
\]
is a forbidden relation, again a contradiction.
\end{itemize}

Therefore at least three elements of $T_a$ are omitted.

\medskip
\noindent\textbf{Lemma 5.}
If $N/12<a\le N/6$, then any admissible set $A\subseteq \{1,\ldots,N\}$ omits at least one element of $T_a$.

\medskip
\noindent\textbf{Proof.}
In this range we still have $2a,3a,6a\le N$. If all three belonged to $A$, then
\[
\frac1{2a}=\frac1{3a}+\frac1{6a}
\]
would be a forbidden relation. Hence at least one of $2a,3a,6a$ is missing, so at least one element of $T_a$ is missing.

\medskip
\noindent\textbf{Step 5: global counting.}

Let $A\subseteq\{1,\ldots,N\}$ be admissible. Because the blocks $T_a$ are pairwise disjoint, omissions forced by different $a\in\mathcal I$ are distinct.

Therefore the number of omitted elements is at least
\[
3I(N/24)+2\bigl(I(N/12)-I(N/24)\bigr)+\bigl(I(N/6)-I(N/12)\bigr).
\]
This simplifies to
\[
I(N/24)+I(N/12)+I(N/6).
\]
Using Lemma~2 three times,
\[
I(N/24)+I(N/12)+I(N/6)
=
\frac25\left(\frac{N}{24}+\frac{N}{12}+\frac{N}{6}\right)+o(N)
=
\frac{7}{60}N+o(N).
\]
Hence
\[
|A|\le N-\frac{7}{60}N+o(N)
=
\left(\frac{53}{60}+o(1)\right)N.
\]

So we obtain the improved upper bound

\medskip
\noindent\textbf{Theorem 6 (Round-2 improvement).}
\[
f(N)\le \left(\frac{53}{60}+o(1)\right)N.
\]

\medskip
\noindent\textbf{Step 6: extending the exact computation.}

Round~1 computed $f(N)$ exactly for $N\le 15$. I extended this to $N\le 26$ by exhaustive search over all subsets of $\{1,\ldots,N\}$, with forbidden hyperedges given by all exact identities
\[
\frac1a=\frac1{b_1}+\cdots+\frac1{b_k}
\]
among distinct integers in $\{1,\ldots,N\}$.

The resulting exact values are
\[
\begin{array}{c|cccccccccccccccccccccccccc}
N
&1&2&3&4&5&6&7&8&9&10&11&12&13&14&15&16&17&18&19&20&21&22&23&24&25&26\\
\hline
f(N)
&1&2&3&4&5&5&6&7&8&9&10&10&11&12&12&13&14&15&16&16&17&18&19&19&20&21
\end{array}
\]

\subsection*{6) ADVERSARIAL VERIFICATION}

\begin{itemize}
\item \emph{Disjointness check.} The residue-pair argument determines the multiplier uniquely, so different blocks really are disjoint.
\item \emph{Local omission check.} Lemma~4 is the delicate point. The proof splits into the only two possible cases: $8a\in A$ or $8a\notin A$. Both are fully closed.
\item \emph{Counting check.} Because blocks are disjoint, local omissions from different blocks cannot overlap. This is exactly what makes the global sum legitimate.
\item \emph{Arithmetic check.} The omitted-density constant is
\[
\frac25\left(\frac1{24}+\frac1{12}+\frac1{6}\right)=\frac25\cdot\frac7{24}=\frac7{60},
\]
so the surviving-density constant is indeed
\[
1-\frac7{60}=\frac{53}{60}.
\]
\item \emph{Compatibility with Round~1.} Since
\[
\frac{53}{60}=0.88333\ldots < 0.892857\ldots=\frac{25}{28},
\]
the new bound is a genuine improvement over the Round-1 result.
\end{itemize}

\subsection*{7) FINAL (EXACTLY ONE)}

\textbf{UNRESOLVED (BUT STRICTLY ADVANCED).}

Round~2 proves the stronger upper bound
\[
f(N)\le \left(\frac{53}{60}+o(1)\right)N,
\]
improving the Round-1 constant $\frac{25}{28}$, and extends the exact computation of $f(N)$ from $N\le 15$ to $N\le 26$.

The original Erd\H{o}s problem remains open, because there is still a large gap between the lower bound $f(N)\ge N/2$ and the new upper bound $f(N)\le (53/60+o(1))N$.

\subsection*{8) COMPLETION ESTIMATE (MANDATORY)}

\textbf{COMPLETION: 58\%}

\subsection*{9) REFERENCES}

\begin{enumerate}
\item T.~F.~Bloom, Erd\H{o}s Problem \#301, \texttt{erdosproblems.com/301}, accessed 2026-03-07.
\item The Round-1 writeup for the previously reconstructed $25/28$ argument of van Doorn.
\item The exact values for $f(N)$ up to $26$ are from the exhaustive computation carried out in this round.
\end{enumerate}
